\begin{abstract}
Speculative decoding improves large language model (LLM) inference throughput by drafting multiple candidate tokens with a smaller model and verifying them with a larger target model. However, reported gains vary substantially across hardware, model pairing, and decoding regime, and consumer-GPU evidence remains limited. We present a controlled empirical study on a single RTX~5090 using a Qwen2.5-3B-Instruct target and same-family 0.5B/1.5B drafts over GSM8K, a 5-subject MMLU subset, and CNN/DailyMail, with frozen manifests and matched prompts across all runs. We evaluate 12 speculative configurations ($k\in\{4,8,16\}$, deterministic/stochastic regimes, two draft sizes) plus regime-matched baselines, and report wall-clock latency, throughput, acceptance dynamics ($\alpha$, $B_{\text{eff}}$), and task-level quality deltas. The best observed configuration is 1.5B, $k{=}8$, deterministic, with $S{=}1.7289$, $\alpha{=}0.3829$, and $B_{\text{eff}}{=}2.96$. Across the grid, performance is non-monotonic in $k$: $k{=}8$ outperforms both $k{=}4$ and $k{=}16$ on mean speedup, indicating an acceptance--cost tradeoff rather than a simple larger-block-is-better effect. Deterministic decoding is consistently faster than stochastic decoding for each draft--$k$ pair, while quality drift is modest on MMLU and CNN/DailyMail but more variable on GSM8K under stochastic settings. We conclude with practical operating recommendations for consumer-GPU deployment and a targeted extension plan to strengthen statistical confidence under fixed additional runtime budget.
\end{abstract}